%% Introduction %% 

Sous le terme de \enquote{chercheurs} nous réunissons les personnes exploitant les collections des institutions patrimoniales à des fins de projets de recherche. 

\subsection{L'accessibilité commence par l'acceptation de la pertinence de la recherche}

%% Situer dans le temps %%

Ce sont les premiers à investir les salles de lecture. Le \gls{td} du dépôt légal est principalement à destination de ce public. C'est le premier à investir l'INAthèque et c'est toujours aujourd'hui vers lui que sont destinés ses services ainsi que les compétences des documentalistes. 

Cela dit, si l'étude des sources \enquote{physiques} : manuscrits, cartes, lettres n'ont plus à prouver leur valeur. Les archives audiovisuelles, à la création de l'INAthèque ne disposaient pas de la même légitimité, ce que déplorait Jean-Noël Jeanneney dans \textit{Le Monde} en 1982 : 
\begin{displayquote}
	Il est temps, en effet, que l'évidence s'en impose à tous, bousculant les timidités intellectuelles et les paresses de l'habitude : la télévision a conquis désormais, dans notre vie collective, une importance trop grande pour que l'apport en soit négligé par l'histoire contemporaine, car la connaissance des trois dernières décennies au moins, en serait amputée d'une dimension cardinale\footcite[][10]{zotero-item-735}.
\end{displayquote}
Aujourd'hui, c'est enfin chose faite, grâce aux avancées législatives, technologiques et sociologiques, l'étude de la télévision et de la radio est enfin démocratisée. Certaines institutions de patrimoines télévisés n'ont pas bénéficié de cette ouvertures à l'étude des documents audiovisuels. En Suisse, par exemple, le dialogue entre les acteurs institutionnels et les historiens n'amènent pas à des solutions durables, d'autant plus que le pays n'a pas été avantagé par sa politique fédéraliste, résistant à la centralisation dont les archives audiovisuelles ont bénéficié en France\footcite[][27]{pradervand2013}. Au Brésil, malgré un dépôt légal, l'absence des documents audiovisuels dans les dispositifs légaux de conservation est l'issue d'un manque de considération de la part des politiques ainsi qu'une absence de reconnaissance de la télévision comme un moyen d'expression culturelle\footcite[][58]{Maria2013}. Face au manque d'engouement pour ces collections, ces pays ne bénéficient ni de collections de diffusions télévisées aussi fournie que celles dont nous disposons avec l'INA et encore moins d'un accès facilité à ces collections. Mais à l'INA, de la même manière que l'étude des documents audiovisuels et en général de la télévision et de la radio ont évolué avec le temps, les chercheurs ont fait de même.

\subsection{Le service de l'INAthèque}

%% Services %%

De tous temps, l'INAthèque et ses documentalistes étaient au service des chercheurs et de leurs usages. Ils ont a leur disposition un outil de consultation des collections nommé \enquote{\gls{hy}} grâce auquel ils peuvent rechercher et consulter les notices. En complément de cet outil, les utilisateurs disposent d'autres outils et services. Cela commence par les connaissances dispensées par les documentalistes qui accompagnent les publics de l'INAthèque à l'utilisation de la base de données, à la compréhension et l'exploitation de ses données. S'ajoute à cela des billets de blog décrivant des fonds récemment traités sur le site \enquote{Hypothèse}\footcite[]{2026} auquel les chercheurs peuvent faire appel pour aiguiller leurs recherches. 

Les chercheurs sont les publics externes à l'INA les plus touchés par ces problématiques de \gls{td} et ils sont demandeurs de discussion traitant de ce sujet. C'est d'autant plus bénéfique qu'étant la cible du \gls{td} du dépôt légal, la réflexion qu'ils peuvent apporter est essentielle. Chez les chercheurs, la tendance est à la collaboration et aux échanges, notamment sur la façon qu'ils ont de chercher les documents et en conséquent sur le \gls{td}. 

C'est également une composante que l'on retrouve avec le Lab, crée en 2022. Précédé par des \enquote{ateliers de recherche méthodologique sur le dépôt légal du web}\footcite[][200]{bermès2024}, des ateliers accompagnant les chercheurs aux archives du web en ligne\footcite[]{webcorpora2017}. Le lab est présenté par l'INA comme un \enquote{incubateur de projets}\footcite{zotero-item-700}. Copiloté par la direction des patrimoines et la direction Data \& Technologies, il est composé d'ingénieurs de la recherche et est recentré sur les données. Il propose le même accompagnement de compréhension et d'analyse des documents de la base de données que les documentalistes en salle de lecture mais il se concentre plus particulièrement sur de plus larges jeux de données en accord avec le sujet de recherche des chercheurs accompagnés. Ils élaborent avec eux leur corpus d'étude, les accompagne sur leur pratique des humanités numériques à partir de données audiovisuelles et sur l'utilisation des outils de l'INA. Cet structure fut crée suite aux demandes des chercheurs de plus en plus récurrentes à réaliser des études quantitatives et qualitatives sur des corpus de données très larges, avec l'assistance d'outils de traitement automatique du langage par exemple\renvoientretien{\AH}{les outils et bases de données de l'INA sur les profils de recherche}{ann:transcriptions}

Loin d'être une initiative propre à l'INA, les \enquote{labs} sont des structures mises en place par différentes institutions patrimoniales afin de former le public à ces nouvelles collections entièrement ou majoritairement numériques et aux outils, méthodologies et corpus qui en découlent\footcite[][196]{bermès2024}. Différentes bibliothèques nationales ont suivi ce mouvement au Royaume-Uni (British Library Labs), aux Pays-Bas (KB Lab), aux Etats-Unis (LC Labs)\footcite[197]{bermès2024} et enfin la BNF avec le BNF DataLab\footcite[][199]{bermès2024}. Toutes ces évolutions du \gls{td} ne nécessitent pas seulement une adaptation des bases de données et des outils de consultation mais également une adaptation des services proposés aux chercheurs afin de garantir leur accessibilité aux collections. D'autant plus que leurs recherches, elles aussi évoluent.

%% Outils %%

Concernant les outils à leur disposition, ils se divisent en outils \enquote{experts} et les outils \enquote{autonomes}. Les outils experts comprennent l'outil de consultation \enquote{\gls{hy}}, permettant de faire des paniers de documents audiovisuels consultables dans \enquote{Mediascope}, qui permet de les segmenter et de les annoter ainsi que dans \enquote{Médiacorpus}, qui comme son nom l'indique permet de créer des corpus. Ces outils sont utilisés à l'INAthèque de la BNF ainsi que dans les six délégations de l'INA situées à Rennes, Toulouse, Strasbourg, Lyon, Marseille et Lille. Ils nécessitent l'accompagnement d'un.e documentaliste, au contraire des outils autonomes, présents dans une cinquantaine de lieux partenaires comme des bibliothèques, cinémathèques ou des archives départementales. Les chercheurs de ces lieux partenaires peuvent accéder au fonds du dépôt légal par une interface nommée \enquote{WebMedia} et \enquote{Opsismédia}, sur lesquels ils peuvent visualiser ces documents. Ces outils ont été construits pour répondre aux besoins des chercheurs mais ils sont aujourd'hui vieillissants, fonctionnant sur des OS peu utilisés et manquant de fonctionnalités pour répondre à de nouveaux besoins tel que la possibilité de créer un compte personnel pour créer et stocker des corpus, à partir desquels ils pourraient créer des datavisualisations. L'INA se place cela dit dans la continuité de son adaptation permanente aux nouveaux besoins de ses utilisateurs et a en développement un nouvel outil pour répondre à ces nouveaux besoins nommé \enquote{Skopus}.\footnote{(Ces outils ont été présenté par Arthur Hénaff lors de son entretien, voir sa partie sur les outils et bases de données de l'INA~\ref{ann:transcriptions},\nameref{ann:transcriptions}), p.~\pageref{ann:transcriptions}.}. 

\subsection{Habitudes de consultation des collections}

Ainsi, nous pouvons diviser les chercheurs de l'INAthèque en deux catégories : les chercheurs du Lab qui sont des utilisateurs réguliers des services de l'INAthèque et les chercheurs plus ponctuels, venant pour des recherches personnelles ou accompagnés ou sur la demande de leur professeur durant leur cursus d'études supérieures. Ils effectuent des recherches sur un sujet restreint, parfois donné par leur professeur et utilisent les outils expert avec l'aide des documentalistes. À ces étudiants, les documentalistes ne leur inculque que ce dont ils ont besoin pour effectuer leur recherche. Selon la période, le thème ou le sujet qu'ils souhaitent étudier, ils leur conseillent d'utiliser plutôt les champs titre, mots-clé, résumé ou même d'utiliser du tout-texte, donc de rechercher dans tous les champs, au risque de générer beaucoup de bruit. Leur sujet de recherche peut également avoir pour objectif de créer des statistiques, ils peuvent dans cet objectif rechercher des informations de dates, de temps et d'heure de diffusion. Dans ce même objectif de réalisation de statistiques, ces chercheurs sont très demandeurs de transcription des émissions afin de réaliser des statistiques sur la fréquence de l'utilisation d'un mot. Mais peu importe leur utilisation des collections, la médiation que leur fournie les documentalistes est indispensable à l'utilisation des outils de l'INAthèque. 

%% Conclusion %% 

La première chose que nous pouvons noter est que bien que le profil des \enquote{chercheurs} est très courant dans les institutions patrimoniales, les chercheurs fréquentant l'INAthèque ont bénéficié de politiques favorables à la conservation des diffusions audiovisuelles. En conséquence, l'INA fait posture de précurseur dans tous les outils qu'elle met en place pour la communication de fonds audiovisuels étant donné que la richesse des collections de cette institution la placent en exception dans le monde patrimonial à l'échelle mondiale. Cela dit, les chercheurs sont les publics les plus touchés par ces enjeux d'accessibilité aux collections des institutions patrimoniales. Ils sont pleinement dépendant des outils crées par les institutions patrimoniales et plus particulièrement de l'expertise des documentalistes. Ces derniers leurs fournissent les clefs pour comprendre et exploiter les outils et collections de l'institution. 